\documentclass[11pt,a4paper]{article}

\usepackage[utf8]{inputenc}
\usepackage[T1]{fontenc}
\usepackage{amsmath,amssymb}
\usepackage{mhchem}
\usepackage{geometry}
\geometry{margin=2.5cm}
\usepackage{hyperref}
\usepackage{enumitem}
\usepackage{xcolor}
\hypersetup{colorlinks=true,linkcolor=black,urlcolor=blue}

\newcommand{\co}{\ensuremath{\mathrm{CO_2}}}
\newcommand{\FE}{\mathrm{FE}}
\newcommand{\kgi}{\mathrm{kg^{-1}}}

\title{\textbf{co2dash --- Methods, Equations, and References}}
\author{}
\date{}

\begin{document}
\maketitle
\vspace{-2.5em}

\noindent Every equation the engine actually evaluates, with the source it is
taken from and the reason for the choice. Equations are transcribed from the
source code (\texttt{techno\_economic.py}, \texttt{proxy.py},
\texttt{calibration.py}, \texttt{uncertainty.py}) so this document matches what
the app computes, not an idealised version.

\begin{quote}
\small\textbf{Citation note.} The references below are the canonical/primary
sources for each method. Exact bibliographic details (year, volume, DOI) should be
re-verified against the original papers before use in a manuscript; they are given
here to document provenance and rationale, not as camera-ready citations.
\end{quote}

\noindent\textbf{Conventions.} All quantities are per kg of product unless
stated. Constants: Faraday $F = 96\,485~\mathrm{C\,mol^{-1}}$;
$1~\mathrm{kWh} = 3.6\times10^{6}~\mathrm{J}$.

\hrulefill

\section{Specific electrical energy}
\emph{Code: }\texttt{specific\_electricity\_kwh\_per\_kg}.
\begin{equation}
E_\mathrm{el} = \frac{n\,F\,V_\mathrm{cell}}{M \cdot \FE \cdot \eta_\mathrm{rect}}
\quad[\mathrm{J\,\kgi}] \;\longrightarrow\; \mathrm{kWh\,\kgi}
\end{equation}
with $n$ = electrons per molecule, $V_\mathrm{cell}$ = full-cell voltage, $M$ =
molar mass of product, $\FE$ = faradaic efficiency, $\eta_\mathrm{rect}$ =
rectifier (AC$\to$DC) efficiency (default $0.95$).

\textbf{Source.} Faraday's law of electrolysis (standard electrochemistry, e.g.
Bard \& Faulkner, \emph{Electrochemical Methods}). The same expression is the
energy core of Jouny, Luc \& Jiao, \emph{Ind.\ Eng.\ Chem.\ Res.} 2018 --- the
study we validate against.

\textbf{Rationale.} This is the single term with the sharpest sensitivity to the
two levers research can move ($\FE$ enters as $1/\FE$; $V_\mathrm{cell}$ enters
linearly). Exposing it explicitly, rather than burying it in an aggregate cost, is
a core design goal. $\eta_\mathrm{rect}$ captures balance-of-plant AC$\to$DC
losses that separate cell energy from plant energy.

\section{Capital recovery factor (levelisation)}
\emph{Code: }\texttt{capital\_recovery\_factor}.
\begin{equation}
\mathrm{CRF} = \frac{i(1+i)^{L}}{(1+i)^{L}-1}, \qquad
\mathrm{CRF} = \tfrac{1}{L}\ \text{if } i=0,
\end{equation}
with $i$ = discount rate, $L$ = plant lifetime (years).

\textbf{Source.} Standard engineering economics (annualised capital charge). The
same device --- an annualised capital charge rate (ACCR $\equiv$ CRF) --- is used
by Osorio-Tejada et al., \emph{Energy Environ.\ Sci.} 2024, our second validation
anchor.

\textbf{Rationale.} We deliberately levelise capital with a CRF rather than run a
full discounted cash-flow with tax and depreciation. This makes the output a
\emph{cost of production} (comparable across routes at a glance) rather than an
after-tax NPV. It is a scope choice, stated openly: absolute LCOP is \emph{not}
expected to equal a levered NPV such as Jouny 2018's, but \emph{is} like-for-like
with Osorio-Tejada 2024 (which uses ACCR).

\section{Levelised cost of product (LCOP)}
\emph{Code: }\texttt{lcop}.
\begin{equation}
\mathrm{LCOP} =
\underbrace{\frac{\mathrm{CAPEX}\cdot\mathrm{CRF} + \mathrm{OPEX}_\mathrm{fix}}
{\dot m_\mathrm{annual}}}_{\text{fixed}}
+ \underbrace{c_{\co}\,m_{\co} + c_\mathrm{el}\,E_\mathrm{el}
+ c_{\mathrm{H_2}}\,m_{\mathrm{H_2}}}_{\text{variable}}
\quad[\$\,\kgi]
\end{equation}

\textbf{Source.} Standard levelised-cost formulation; structure follows the TEA
breakdowns in Jouny 2018 and Osorio-Tejada 2024 (feedstock + electricity +
annualised capital + fixed O\&M).

\textbf{Rationale.} The variable terms are physical and largely model-independent
(feedstock \co, electricity via $E_\mathrm{el}$, optional \ce{H2} co-feed), so
they reproduce published numbers to within a few percent. The fixed term carries
the model's largest uncertainty because CAPEX is an estimate (\S8).

\section{Net \texorpdfstring{\co}{CO2} abatement (life-cycle / climate)}
\emph{Code: }\texttt{net\_abatement\_kg\_per\_kg}.
\begin{equation}
\text{net} = (1-\varphi)\,m_{\co}
\;-\;\bigl(I_\mathrm{grid}\,E_\mathrm{el} + e_\mathrm{capture} + e_\mathrm{process}\bigr)
\end{equation}
with $\varphi$ = end-of-life release fraction ($0$ = durable product, $1$ = fuel
re-oxidised), $I_\mathrm{grid}$ = grid carbon intensity
($\mathrm{kg\,\co\,kWh^{-1}}$), $e_\mathrm{capture}, e_\mathrm{process}$ = \co\
penalties of capture and other process steps.

\textbf{Source.} Life-cycle CCU accounting; the utilisation-vs-sequestration
distinction follows CCU LCA methodology (e.g.\ von der Assen, Bardow and
co-workers) and IPCC-style grid-emission accounting.

\textbf{Rationale.} The $(1-\varphi)$ term is the crux of \emph{utilisation} vs
\emph{storage}: a fuel later burned ($\varphi\to 1$) keeps almost none of the \co\
credit. Making $\varphi$ an explicit lever prevents counting utilised \co\ as
permanently abated. Grid intensity multiplies the (large) electrical energy, which
is why the climate verdict is dominated by the electricity source.

\section{Break-even grid intensity}
\emph{Code: }\texttt{breakeven\_grid\_intensity} (net $=0$ solved for
$I_\mathrm{grid}$).
\begin{equation}
I^{*} = \frac{(1-\varphi)\,m_{\co} - e_\mathrm{capture} - e_\mathrm{process}}
{E_\mathrm{el}}
\end{equation}
Above $I^{*}$, the conversion \emph{increases} emissions.

\textbf{Source.} Algebraic rearrangement of Section~4 (no external source).

\textbf{Rationale.} Given in closed form and surfaced as a first-class KPI because
``how clean must my electricity be for this to help at all?'' is one of the most
decision-relevant questions, and it has an exact answer.

\section{Marginal abatement cost (MAC)}
\emph{Code: }\texttt{marginal\_abatement\_cost}.
\begin{equation}
\mathrm{MAC} = \frac{\mathrm{LCOP}_\mathrm{CCU} - \mathrm{LCOP}_\mathrm{conv}}
{\text{net}}\quad[\$\,\kgi\,\co], \qquad
\mathrm{MAC} = +\infty\ \text{if net} \le 0,
\end{equation}
($\times 1000$ for $\$\,\mathrm{tonne^{-1}}$).

\textbf{Source.} Standard definition of marginal abatement cost from
climate/energy economics (cost premium over the incumbent per unit \co\ avoided).

\textbf{Rationale.} MAC is the single figure investors and regulators read, and it
couples cost and climate into one comparable number. Returning $+\infty$ when
net $\le 0$ is deliberate: a route that is not climate-positive has no defined
abatement cost, and hiding that behind a finite number would mislead.

\section{Activity proxy --- Computational Hydrogen Electrode (CHE)}
\emph{Code: }\texttt{proxy.py}. Turns DFT intermediate energies into an activity
target when experimental $\FE$ is unavailable.

Formation free energy of each adsorbed intermediate (\co\ reference $=0$):
\begin{equation}
G_f({}^{*}X) = \Delta E_f({}^{*}X) + \Delta(\mathrm{ZPE}-T\Delta S)_X .
\end{equation}
Per proton-coupled electron-transfer (PCET) step, and the limiting potential:
\begin{equation}
\Delta G_i(U) = \Delta G_i(0) - eU, \qquad
U_L = -\max_i \Delta G_i(0)\quad (e = 1\ \text{in eV/V}).
\end{equation}
Overpotential and the transparent activity$\to$voltage map fed to the TEA:
\begin{equation}
\eta = U_\mathrm{eq} - U_L, \qquad
V_\mathrm{cell} \approx V_\mathrm{baseline} + \eta .
\end{equation}
Equilibrium potentials used (V vs RHE): CO $-0.10$, HCOOH $-0.12$,
\ce{CH3OH} $+0.02$.

\textbf{Source.} N{\o}rskov et al., \emph{J.\ Phys.\ Chem.\ B} 2004 (computational
hydrogen electrode); \co RR pathway/limiting-potential analysis following
Peterson \& N{\o}rskov and the \co-reduction scaling-relation literature.

\textbf{Rationale.} CHE is the standard, transparent way to rank catalyst
\emph{activity} from adsorption energies without kinetics. We map activity to
\textbf{cell voltage} (which the TEA already consumes), not to $\FE$, because
activity $\neq$ selectivity --- predicting $\FE$ from thermodynamic descriptors is
not defensible, so we do not claim it. \emph{Honesty defaults:}
$\Delta(\mathrm{ZPE}-T\Delta S)$ corrections are \textbf{off} by default
(electronic energies only); equilibrium potentials are textbook values to verify
for pH/reference convention; $V_\mathrm{baseline}$ is an explicit adjustable
parameter, not a first-principles cell model.

\section{Uncertainty propagation (Monte-Carlo)}
\emph{Code: }\texttt{uncertainty.py::propagate\_mc}. Each uncertain input is
sampled from its distribution and pushed through the deterministic engine
(Sections~1--6) $N$ times (default $\approx 4\times 10^{4}$):
\begin{itemize}[nosep]
  \item Normal: $x\sim\mathcal N(\mu,\sigma)$, truncated to physical bounds.
  \item Lognormal: $x = a\,e^{\mathcal N(0,\ln b)}$ (median $a$, geometric std
  $b$) for strictly-positive, right-skewed quantities.
\end{itemize}
Outputs are read empirically from the sample:
\begin{equation}
P(\text{net}>0) = \frac1N\sum_k \mathbf 1[\text{net}_k>0], \qquad
P(\mathrm{MAC}<c) = \frac1N\sum_k \mathbf 1[\mathrm{MAC}_k<c],
\end{equation}
with P05/P95 percentiles reported for the MAC distribution. The per-field
\emph{data tier} sets the noise floor $\sigma$: \textsc{computed} $\approx 5\%$,
\textsc{lab\_validated} $\approx 3\%$, \textsc{lit\_extracted} $\approx 20\%$,
\textsc{estimated} $\approx 40\%$ (typically CAPEX).

\textbf{Source.} Monte-Carlo propagation of distributions (GUM Supplement~1, JCGM
101:2008).

\textbf{Rationale.} The engine is deterministic; all uncertainty is in the inputs,
so propagating input distributions is the correct, assumption-light way to obtain
honest output distributions. Probabilities are empirical frequencies, not
parametric approximations. Tiers tie the noise to data provenance, so the width of
the answer reflects how well the inputs are known.

\section{Uncertainty calibration}
\emph{Code: }\texttt{calibration.py}. Applied to a learned surrogate's predictive
std \emph{before} it enters Section~8.

Empirical coverage of central Gaussian intervals (diagnosis):
\begin{equation}
\widehat{\mathrm{cov}}(\ell) = \frac1m\sum_i
\mathbf 1\!\left[\,|y_i-\mu_i| \le z_\ell\,\sigma_i\,\right], \qquad
z_\ell = \Phi^{-1}\!\left(\tfrac12 + \tfrac{\ell}{2}\right).
\end{equation}
Temperature scaling (global std correction), fitted by Gaussian MLE:
\begin{equation}
s^2 = \frac1m\sum_i \frac{(y_i-\mu_i)^2}{\sigma_i^2}, \qquad
\sigma_i' = s\,\sigma_i .
\end{equation}
Normalised split-conformal prediction (distribution-free, finite-sample):
\begin{equation}
r_i = \frac{|y_i-\mu_i|}{\sigma_i}, \qquad
q = \big\lceil (m+1)(1-\alpha)\big\rceil / m\ \text{empirical quantile of }\{r_i\}, \qquad
\text{interval} = \mu \pm q\,\sigma .
\end{equation}

\textbf{Source.} Calibrated regression: Kuleshov, Fenner \& Ermon, \emph{ICML}
2018; Levi et al.\ 2022; the scalar temperature idea generalises Guo et al.,
\emph{ICML} 2017. Split/normalised conformal: Vovk, Gammerman \& Shafer 2005; Lei
et al., \emph{JASA} 2018; Papadopoulos et al.\ 2008.

\textbf{Rationale.} An over-confident surrogate yields feasibility probabilities
that are confident \emph{and wrong} --- the worst failure for a decision tool.
Temperature scaling is a one-parameter, robust fix for systematic
over/under-confidence; split-conformal adds a distribution-free finite-sample
coverage guarantee. The conformal half-width is mapped back to a
Gaussian-equivalent std so the existing Monte-Carlo sampler consumes calibrated
uncertainty unchanged (a documented width-matching approximation).

\section{Global sensitivity (Sobol)}
\emph{Code: }\texttt{uncertainty.py::sobol\_indices} (via SALib). Variance-based
total-effect indices $S_T$ are computed with Saltelli sampling over the input
ranges, ranking each input by its share of MAC variance.

\textbf{Source.} Sobol, \emph{Math.\ Comput.\ Simul.} 2001; Saltelli et al.\ 2010
(estimators SALib implements).

\textbf{Rationale.} Variance-based (rather than local/one-at-a-time) sensitivity
suits a nonlinear model with inputs uncertain over wide ranges; total-effect
indices capture interactions. The tornado of $S_T$ tells the user \emph{what to
work on} (catalyst vs energy vs plant).

\section*{Consolidated references}
\begin{enumerate}[nosep,leftmargin=1.4em]
\item Bard, A.\,J.; Faulkner, L.\,R. \emph{Electrochemical Methods.} (Faraday's law.)
\item Jouny, M.; Luc, W.; Jiao, F. \emph{A General Techno-Economic Analysis of
\co\ Electrolysis Systems.} Ind.\ Eng.\ Chem.\ Res., 2018. (TEA anchor.)
\item Osorio-Tejada, J. et al. \emph{Techno-economic analysis of \co\ conversion.}
Energy Environ.\ Sci., 2024. (Like-for-like levelised-cost anchor; ACCR $=$ CRF.)
\item von der Assen, N.; Bardow, A. et al. --- \co-utilisation LCA methodology.
\item JCGM 101:2008 --- \emph{Supplement 1 to the GUM} (Monte-Carlo propagation).
\item N{\o}rskov, J.\,K. et al.\ --- computational hydrogen electrode, J.\ Phys.\
Chem.\ B, 2004.
\item Peterson, A.\,A.; N{\o}rskov, J.\,K. --- \co-reduction limiting potential /
scaling relations.
\item Kuleshov, V.; Fenner, N.; Ermon, S. \emph{Accurate Uncertainties for Deep
Learning Using Calibrated Regression.} ICML, 2018.
\item Levi, D. et al. \emph{Evaluating and Calibrating Uncertainty Prediction in
Regression Tasks.} 2022.
\item Guo, C. et al. \emph{On Calibration of Modern Neural Networks.} ICML, 2017.
\item Vovk, V.; Gammerman, A.; Shafer, G. 2005; Lei, J. et al. \emph{JASA} 2018;
Papadopoulos, H. et al.\ 2008 (conformal prediction).
\item Sobol, I.\,M. \emph{Math.\ Comput.\ Simul.} 2001; Saltelli, A. et al.\ 2010
(variance-based sensitivity; SALib).
\end{enumerate}

\noindent\small Grid carbon intensities used in scenarios: Ember / IEA 2024;
Australian National Greenhouse Accounts Factors 2025. Dataset citations for the
ML/calibration work are in \texttt{docs/VALIDATION.md}.

\end{document}
